% Unit 091 terjemahan Bahasa Indonesia dari chapter7.tex baris 712--891.
% Sumber: Methods of Algebra, Volume 2, commit
% 9a5803ff77dd3257484cb177f851a73770a59dd3 (CC BY 4.0).
% stable-unit-id: o014.aljabr2.chapter7.closed-structure-dgcat
% Cakupan: Bagian 7.4 lengkap, ``Studi Kasus: Struktur Tertutup pada
% Kategori-dg''; berhenti tepat sebelum Bagian 7.5 pada baris 892.

% segment-id: o014.li2.chapter7.closed-structure-dgcat.h001
\section{Studi Kasus: Struktur Tertutup pada Kategori-dg}\label{sec:dg-closed}
\hypertarget{o014.aljabr2.chapter7.closed-structure-dgcat}{ }
% segment-id: o014.li2.chapter7.closed-structure-dgcat.p001
Sebelumnya telah dibahas struktur kategori monoidal tertutup pada
$\cate{Cat}$, dengan $\otimes$ berasal dari produk kategori biasa
$\mathcal{C}_1\times\mathcal{C}_2$. Di sisi lain, relatif terhadap gelanggang
komutatif $\Bbbk$ yang dipilih, Definisi~\ref{def:dgCat} memperkenalkan
kategori-dg, yakni kategori yang diperkaya atas
$(\cate{C}(\Bbbk\dcate{Mod}),\otimes)$. Bagian ini bertujuan menunjukkan bahwa
semua kategori-dg kecil juga membentuk kategori monoidal tertutup
$\cate{dgCat}_{\Bbbk}$. Kategori ini dapat dipandang sebagai versi
linearisasi sekaligus versi kompleks dari $\cate{Cat}$ dan muncul secara
alami serta berguna dalam banyak keadaan. Yang diperlukan hanyalah beberapa
konstruksi yang panjang tetapi tidak sulit.

% segment-id: o014.li2.chapter7.closed-structure-dgcat.p002
Langkah pertama ialah mengangkat produk kategori biasa menjadi produk
tensor kategori-dg. Hal ini melibatkan struktur anyaman simetris pada
$(\cate{C}(\Bbbk\dcate{Mod}),\otimes)$.

% segment-id: o014.li2.chapter7.closed-structure-dgcat.d001
\begin{definition}\label{def:dgCat-tensor}
	Misalkan $\mathcal{C}_1$ dan $\mathcal{C}_2$ kategori-dg di atas $\Bbbk$.
	Definisikan kategori-dg baru $\mathcal{C}_1\otimes\mathcal{C}_2$ dengan
	% segment-id: o014.li2.chapter7.closed-structure-dgcat.q001
	\begin{align*}
		\Obj(\mathcal{C}_1 \otimes \mathcal{C}_2) & := \Obj(\mathcal{C}_1) \times \Obj(\mathcal{C}_2), \\
		\Hom_{\mathcal{C}_1 \otimes \mathcal{C}_2}^\bullet((X, X'), (Y, Y')) & := \Hom_{\mathcal{C}_1}^\bullet(X, Y) \otimes \Hom_{\mathcal{C}_2}^\bullet(X', Y'),
	\end{align*}
	dan tetapkan $\identity_X\otimes\identity_{X'}$ sebagai satuan dari
	$\Hom_{\mathcal{C}_1\otimes\mathcal{C}_2}^\bullet((X,X'),(X,X'))$.
	Komposisi morfisme didefinisikan melalui
	% segment-id: o014.li2.chapter7.closed-structure-dgcat.q002
	\[\begin{tikzcd}
		\left( \Hom_{\mathcal{C}_1}^\bullet(Y, Z) \otimes \Hom_{\mathcal{C}_2}^\bullet(Y', Z')\right) \otimes
		\left( \Hom_{\mathcal{C}_1}^\bullet(X, Y) \otimes \Hom_{\mathcal{C}_2}^\bullet(X', Y') \right) \arrow[d, "\sim" sloped] \\
		\left( \Hom_{\mathcal{C}_1}^\bullet(Y, Z) \otimes \Hom_{\mathcal{C}_1}^\bullet(X, Y) \right) \otimes
		\left( \Hom_{\mathcal{C}_2}^\bullet(Y', Z') \otimes \Hom_{\mathcal{C}_2}^\bullet(X', Y') \right) \arrow[d] \\
		\Hom_{\mathcal{C}_1}^\bullet(X, Z) \otimes \Hom_{\mathcal{C}_2}^\bullet(X', Z')
	\end{tikzcd}\]
	dengan isomorfisme di dalamnya berasal dari struktur anyaman Koszul dan
	kendala asosiativitas.
\end{definition}

% segment-id: o014.li2.chapter7.closed-structure-dgcat.p003
Dari simetri struktur anyaman dapat diperiksa bahwa
$\mathcal{C}_1\otimes\mathcal{C}_2$ secara alami ekuivalen dengan
$\mathcal{C}_2\otimes\mathcal{C}_1$. Tuliskan $\Bbbk$ untuk kategori-dg yang
hanya mempunyai satu objek dan kompleks endomorfisme $\Bbbk$. Mudah
dibuktikan bahwa $\Bbbk\otimes\mathcal{C}$ ekuivalen dengan
$\mathcal{C}\otimes\Bbbk$.

% segment-id: o014.li2.chapter7.closed-structure-dgcat.d002
\begin{definition}\label{def:dgCat-as-cat}
	\index[sym1]{dgCat@$\cate{dgCat}_{\Bbbk}$}
	Misalkan $\Bbbk$ gelanggang komutatif. Tuliskan $\cate{dgCat}_{\Bbbk}$
	untuk kategori semua kategori-dg kecil di atas $\Bbbk$. Kategori ini adalah
	kategori monoidal simetris dengan satuan $\Bbbk$.
\end{definition}

% segment-id: o014.li2.chapter7.closed-structure-dgcat.p004
Pembatasan pada kategori kecil tentu dimaksudkan untuk menghindari persoalan
teori himpunan. Manfaat utamanya ialah bahwa semua funktor antara dua kategori
kecil $\mathcal{C}$ dan $\mathcal{D}$ yang diberikan membentuk himpunan kecil.

% segment-id: o014.li2.chapter7.closed-structure-dgcat.p005
Langkah kedua ialah memperkaya himpunan $\Hom$ antara funktor-dg menjadi
kompleks.

% segment-id: o014.li2.chapter7.closed-structure-dgcat.d003
\begin{definition}
	Misalkan $F,G:\mathcal{C}\to\mathcal{D}$ funktor-dg antara kategori-dg.
	Untuk setiap $n\in\Z$, suatu morfisme berderajat $n$ dari $F$ ke $G$
	didefinisikan sebagai data
	% segment-id: o014.li2.chapter7.closed-structure-dgcat.q003
	\[ \phi = \left( \phi_X \right)_{X \in \Obj(\mathcal{C})}, \quad \phi_X \in \Hom_{\mathcal{D}}^n(FX, GX), \]
	dengan syarat bahwa untuk semua $m\in\Z$, $X,Y\in\Obj(\mathcal{C})$, dan
	$f\in\Hom^m_{\mathcal{C}}(X,Y)$, berlaku persamaan berikut dalam
	$\Hom^{n+m}_{\mathcal{D}}(FX,GY)$:
	% segment-id: o014.li2.chapter7.closed-structure-dgcat.q004
	\[ (Gf) \phi_X = (-1)^{nm} \phi_Y (Ff); \]
	komposisi dalam rumus ini dipahami sebagai morfisme dalam
	$\Bbbk\dcate{Mod}$
	% segment-id: o014.li2.chapter7.closed-structure-dgcat.q005
	\begin{gather*}
		\Hom^m_{\mathcal{D}}(GX, GY) \otimes \Hom^n_{\mathcal{D}}(FX, GX) \to \Hom^{m+n}_{\mathcal{D}}(FX, GY), \\
		\Hom^n_{\mathcal{D}}(FY, GY) \otimes \Hom^m_{\mathcal{D}}(FX, FY) \to \Hom^{m+n}_{\mathcal{D}}(FX, GY).
	\end{gather*}
	Modul-$\Bbbk$ yang terdiri atas semua morfisme berderajat $n$
	$\phi:F\to G$ ditulis $\iHom^n(F,G)$.
\end{definition}

% segment-id: o014.li2.chapter7.closed-structure-dgcat.p006
Kondisi di atas dapat dipandang sebagai pernyataan bahwa diagram morfisme
berderajat
% segment-id: o014.li2.chapter7.closed-structure-dgcat.q006
\[\begin{tikzcd}
	FX \arrow[r, "{\phi_X}"] \arrow[d, "Ff"'] & GX \arrow[d, "Gf"] \\
	FY \arrow[r, "{\phi_Y}"'] & GY
\end{tikzcd} \quad
\begin{array}{cc}
	\phi: \text{berderajat }n \\
	f: \text{berderajat }m
\end{array}\]
komutatif hingga tanda $(-1)^{nm}$ yang diwajibkan oleh aturan tanda Koszul.
Secara ketat, ``morfisme berderajat'' tersebut bukanlah morfisme, melainkan
harus dipahami sebagai unsur kompleks $\Hom$.

% segment-id: o014.li2.chapter7.closed-structure-dgcat.dp001
\begin{definition-proposition}
	Untuk funktor-dg $F,G:\mathcal{C}\to\mathcal{D}$ dan setiap $n\in\Z$,
	homomorfisme berikut dapat didefinisikan:
	% segment-id: o014.li2.chapter7.closed-structure-dgcat.q007
	\[ d^n = d^n_{\iHom^\bullet(F, G)}: \iHom^n(F, G) \to \iHom^{n+1}(F, G). \]

	Untuk setiap $\phi=(\phi_X)_{X\in\Obj(\mathcal{C})}$, tetapkan
	% segment-id: o014.li2.chapter7.closed-structure-dgcat.q008
	\[ (d^n \phi)_X := d^n_{\Hom^\bullet(FX, GX)}(\phi_X). \]
	Hal ini membuat $\left(\iHom^n(F,G),d^n\right)_{n\in\Z}$ menjadi kompleks
	$\iHom^\bullet(F,G)$.
\end{definition-proposition}
% segment-id: o014.li2.chapter7.closed-structure-dgcat.pf001
\begin{proof}
	Pertama-tama, periksa bahwa $d^n\phi$ memang termasuk
	$\iHom^{n+1}(F,G)$. Misalkan $f\in\Hom^m_{\mathcal{C}}(X,Y)$. Terapkan
	$d^{m+n}$ pada kedua ruas
	$(Gf)\phi_X=(-1)^{nm}\phi_Y(Ff)$ (subskrip
	$\Hom^\bullet(FX,GY)$ dihilangkan); diperoleh
	% segment-id: o014.li2.chapter7.closed-structure-dgcat.q009
	\begin{align*}
		d^{m+n}((Gf) \phi_X) & = (d^m Gf) \phi_X + (-1)^m (Gf) (d^n \phi_X) \\
		& = G (d^m f) \phi_X + (-1)^m (Gf) (d^n \phi)_X \\
		& = (-1)^{n(m+1)} \phi_Y F(d^m f) + (-1)^m (Gf) (d^n \phi)_X, \\
		(-1)^{nm} d^{m+n}(\phi_Y (Ff)) & = (-1)^{nm} (d^n \phi_Y) Ff + (-1)^{nm+n} \phi_Y d^n(Ff) \\
		& = (-1)^{nm} (d^n \phi)_Y Ff + (-1)^{nm+n} \phi_Y F(d^m f).
	\end{align*}
	Dari kesamaan kedua ekspresi itu segera diperoleh
	$(Gf)(d^n\phi)_X=(-1)^{(n+1)m}(d^n\phi)_YFf$.

	Kedua, karena $\Hom^\bullet(FX,GX)$ adalah kompleks,
	% segment-id: o014.li2.chapter7.closed-structure-dgcat.q010
	\begin{align*}
		(d^{n+1} d^n \phi)_X & = d^{n+1}_{\Hom^\bullet(FX, GX)}((d^n \phi)_X) \\
		& = d^{n+1}_{\Hom^\bullet(FX, GX)} d^n_{\Hom^\bullet(FX, GX)} \phi_X = 0.
	\end{align*}
	Jadi, diperoleh kompleks $\iHom^\bullet(F,G)$.
\end{proof}

% segment-id: o014.li2.chapter7.closed-structure-dgcat.p007
Untuk tiga funktor-dg $F,G,H:\mathcal{C}\to\mathcal{D}$, terdapat operasi
komposisi
% segment-id: o014.li2.chapter7.closed-structure-dgcat.q011
\begin{equation}\label{eqn:dgfct-iHom}\begin{aligned}
	\iHom^a(G, H) \otimes \iHom^b(F, G) & \to \iHom^{a+b}(F, H) \\
	\psi \otimes \phi & \mapsto \psi\phi := \left( \psi_X \phi_X \right)_{X \in \Obj(\mathcal{C})},
\end{aligned}\end{equation}
dengan $a,b\in\Z$. Operasi ini juga asosiatif dan memenuhi
$d^{a+b}(\psi\phi)=(d^a\psi)\phi+(-1)^a\psi(d^b\phi)$. Menurut definisi,
	semuanya cukup diperiksa pada
$\Hom^\bullet_{\mathcal{D}}$. Komposisi ini harus dipahami sebagai komposisi
vertikal morfisme berderajat, yang digambarkan sebagai
% segment-id: o014.li2.chapter7.closed-structure-dgcat.q012
\[\begin{tikzcd}
	\mathcal{C}
	\arrow[bend left=70, rr, "F", ""' name=U]
	\arrow[rr, "G" name=MM, ""' name=M]
	\arrow[bend right=70, rr, "" name=D, "H"'] & &
	\arrow[Rightarrow, to path=(U) -- (MM) \tikztonodes, "\phi"] \arrow[Rightarrow, to path=(M) -- (D) \tikztonodes, "\psi"] \mathcal{D}
\end{tikzcd}  \quad \text{dikomposisikan menjadi} \quad \begin{tikzcd}
	\mathcal{C}
	\arrow[bend left=50, rr, "F", ""' name=U]
	\arrow[bend right=50, rr, "" name=D, "H"']
	& & \arrow[Rightarrow, to path=(U) -- (D) \tikztonodes, "\psi \phi"] \mathcal{D} .
\end{tikzcd}\]

% segment-id: o014.li2.chapter7.closed-structure-dgcat.p008
Tinjau kasus khusus $n=0$. Syarat untuk
$\phi=(\phi_X)_X\in\iHom^0(F,G)$ sama artinya dengan
$(Gf)\phi_X=\phi_Y(Ff)$ untuk semua $f\in\Hom^m(X,Y)$, sedangkan syarat
$d^0\phi=0$ sama artinya dengan
$\phi_X\in\Ker\left(d^0_{\Hom^\bullet(FX,GX)}\right)$ untuk setiap
	$X\in\Obj(\mathcal{C})$. Jadi, unsur modul-$\Bbbk$
% segment-id: o014.li2.chapter7.closed-structure-dgcat.q013
\[ \Ker\left[ \iHom^0(F, G) \xrightarrow{d^0} \iHom^1(F, G) \right] \]
tepat merupakan morfisme $\phi:F\to G$ antara funktor-dg $F$ dan $G$ dalam
pengertian klasik.

% segment-id: o014.li2.chapter7.closed-structure-dgcat.p009
Langkah berikutnya ialah memperkaya kategori funktor antara kategori-dg
menjadi versi dg.

% segment-id: o014.li2.chapter7.closed-structure-dgcat.d004
\begin{definition}
	Untuk sembarang kategori-dg $\mathcal{C}$ dan $\mathcal{D}$, definisikan
	kategori-dg $\iHom(\mathcal{C},\mathcal{D})$ sebagai berikut.
	\begin{itemize}
		\item Objek-objeknya adalah semua funktor-dg
		$F:\mathcal{C}\to\mathcal{D}$.
		\item Untuk sembarang funktor-dg $F,G:\mathcal{C}\to\mathcal{D}$,
		kompleks $\Hom$ antara keduanya didefinisikan sebagai
		$\iHom^\bullet(F,G)$ yang telah didefinisikan di atas.
		\item Operasi komposisi
		$\iHom^\bullet(G,H)\otimes\iHom^\bullet(F,G)\to
		\iHom^\bullet(F,H)$ didefinisikan seperti dalam
		\eqref{eqn:dgfct-iHom}; operasi ini disebut komposisi vertikal.
	\end{itemize}
\end{definition}

% segment-id: o014.li2.chapter7.closed-structure-dgcat.p010
Jika $\mathcal{C}$ dan $\mathcal{D}$ keduanya kategori-dg kecil, maka
$\iHom(\mathcal{C},\mathcal{D})$ juga kecil: himpunan objeknya merupakan
himpunan kecil.

% segment-id: o014.li2.chapter7.closed-structure-dgcat.p011
Dalam Definisi~\ref{def:dgCat-tensor}, kita telah mendefinisikan hasil kali tensor
dua kategori-dg, sehingga kategori $\cate{dgCat}_{\Bbbk}$ dari
Definisi~\ref{def:dgCat-as-cat} menjadi kategori monoidal simetris terhadap
$\otimes$. Kini hubungan antara $\otimes$ dan $\iHom$ dapat dinyatakan dengan
jelas sebagai berikut.

% segment-id: o014.li2.chapter7.closed-structure-dgcat.pr001
\begin{proposition}\label{prop:dgCat-closed}
	Kategori monoidal simetris $\cate{dgCat}_{\Bbbk}$ adalah tertutup. Hom
	internalnya diberikan oleh $\iHom(\mathcal{C},\mathcal{D})$ yang
	didefinisikan di atas, dengan $\mathcal{C}$ dan $\mathcal{D}$ merentang
	semua kategori-dg kecil.
\end{proposition}
% segment-id: o014.li2.chapter7.closed-structure-dgcat.pf002
\begin{proof}
	Kuncinya ialah memberikan isomorfisme kanonik
	% segment-id: o014.li2.chapter7.closed-structure-dgcat.q014
	\[ \Hom_{\cate{dgCat}_{\Bbbk}}(\mathcal{C} \otimes \mathcal{D}, \mathcal{E}) \simeq \Hom_{\cate{dgCat}_{\Bbbk}}(\mathcal{C}, \iHom(\mathcal{D}, \mathcal{E}) ). \]

	Jika struktur dg diabaikan, masalahnya mudah: memberikan
	$F:\mathcal{C}\times\mathcal{D}\to\mathcal{E}$ sama artinya dengan
	memberikan funktor $\mathcal{C}\to\mathcal{E}^{\mathcal{D}}$ yang
	memetakan objek $X$ ke funktor
	$F(X,\mathord\cdot):\mathcal{D}\to\mathcal{E}$. Pokok persoalannya ialah
	syarat pada tingkat $\Hom$ agar sebuah funktor menjadi funktor-dg.

	Menurut definisi, membuat $F$ menjadi funktor-dg sama artinya dengan
	mengangkat pemetaan pada himpunan $\Hom$ menjadi morfisme kompleks
	% segment-id: o014.li2.chapter7.closed-structure-dgcat.q015
	\[ F_{X, Y, Z, W}: \Hom^\bullet_{\mathcal{C}}(X, Y) \otimes \Hom^\bullet_{\mathcal{D}}(Z, W) \to \Hom^\bullet_{\mathcal{E}}(F(X, Z), F(Y, W)), \]
	yang dituntut funktorial dalam keempat peubah. Dengan membahas kedua peubah
	dari $F$ secara terpisah, data ini dapat diuraikan menjadi dua bagian:
	\begin{itemize}
		\item untuk setiap $X\in\Obj(\mathcal{C})$, funktor
		$F(X,\mathord\cdot)$ diangkat menjadi funktor-dg;
		\item diberikan keluarga morfisme
		$F_{X,Y,Z}:\Hom^\bullet_{\mathcal{C}}(X,Y)\to
		\Hom^\bullet_{\mathcal{E}}(F(X,Z),F(Y,Z))$ yang funktorial dalam ketiga
		peubah.
	\end{itemize}
	Sebagai ilustrasi, $F_{X,Y,Z,W}(f\otimes g)$ diperoleh dari
	$F_{X,Y,Z}(f)$ dengan menggunakan
	$g\in\Hom^b_{\mathcal{D}}(Z,W)$ melalui $F(Y,\mathord\cdot)$. Data pada butir
	kedua sama artinya dengan keluarga morfisme
	% segment-id: o014.li2.chapter7.closed-structure-dgcat.q016
	\[ \Hom^\bullet_{\mathcal{C}}(X, Y) \to \iHom^\bullet(F(X, \cdot), F(Y, \cdot)), \]
	yang funktorial dalam $X$ dan $Y$. Kesimpulannya, memberikan funktor-dg
$F$ sama artinya dengan memberikan funktor-dg dari $\mathcal{C}$ ke
$\iHom(\mathcal{D},\mathcal{E})$.
\end{proof}

% segment-id: o014.li2.chapter7.closed-structure-dgcat.p012
Bersama teori dalam \S\ref{sec:closed-monoidal}, pembahasan di atas
menunjukkan bahwa $\cate{dgCat}_{\Bbbk}$ diperkaya atas dirinya sendiri:
untuk sembarang kategori-dg kecil $\mathcal{C}$ dan $\mathcal{D}$, objek
$\iHom$ antara keduanya tetap merupakan kategori-dg.

% segment-id: o014.li2.chapter7.closed-structure-dgcat.p013
Karena $\otimes$ dan $\iHom$ saling menentukan melalui adjungsi
\eqref{eqn:closed-monoidal-adj0}, setelah menerima definisi salah satu di
antara $\otimes$ dan $\iHom$, definisi yang lain setidaknya sama wajarnya.

% segment-id: o014.li2.chapter7.closed-structure-dgcat.r001
\begin{remark}
	Dengan Proposisi~\ref{prop:dgCat-closed} di tangan, sifat umum kategori
	monoidal tertutup menunjukkan bahwa untuk tiga kategori-dg kecil sembarang
	$\mathcal{C}$, $\mathcal{D}$, dan $\mathcal{E}$, komposisi funktor dapat
	diangkat menjadi funktor-dg
	% segment-id: o014.li2.chapter7.closed-structure-dgcat.q017
	\[ \iHom(\mathcal{D}, \mathcal{E}) \otimes \iHom(\mathcal{C}, \mathcal{D}) \to \iHom(\mathcal{C}, \mathcal{E}). \]
	Pada tingkat objek, funktor tersebut memetakan $(G,F)$ ke $GF$; pada
	tingkat kompleks $\Hom$, funktor itu ditulis
	% segment-id: o014.li2.chapter7.closed-structure-dgcat.q018
	\begin{equation}\begin{gathered}\label{eqn:dgCat-fct-compose}
		\begin{tikzcd}[row sep=small, column sep=small]
			\iHom^\bullet(G_1, G_2) \otimes \iHom^\bullet(F_1, F_2) \arrow[r] \arrow[equal, d] & \iHom^\bullet(G_1 F_1, G_2 F_2) \\
			\Hom^\bullet_{ \iHom(\mathcal{D}, \mathcal{E}) \otimes \iHom(\mathcal{C}, \mathcal{D}) }((G_1, F_1), (G_2, F_2)) &
		\end{tikzcd} \\
		F_i: \mathcal{C} \to \mathcal{D}, \quad G_i: \mathcal{D} \to \mathcal{E}, \quad i = 1, 2 .
	\end{gathered}\end{equation}

	Operasi di atas dapat dibayangkan sebagai komposisi horizontal morfisme
	antara funktor-dg, dengan diagram
	% segment-id: o014.li2.chapter7.closed-structure-dgcat.q019
	\[ \begin{tikzcd}
		\mathcal{C} \arrow[bend left=50, r, "F_1", ""' name=LU] \arrow[bend right=50, r, "" name=LD, "F_2"'] &
		\mathcal{D} \arrow[bend left=50, r, "G_1", ""' name=RU] \arrow[bend right=50, r, "" name=RD, "G_2"'] &
		\arrow[Rightarrow, to path=(LU) -- (LD) \tikztonodes, "\phi"] \arrow[Rightarrow, to path=(RU) -- (RD) \tikztonodes, "\psi"] \mathcal{E}
	\end{tikzcd} \quad \text{dikomposisikan menjadi} \quad \begin{tikzcd}
		\mathcal{C}
		\arrow[bend left=50, rr, "G_1 F_1", ""' name=U]
		\arrow[bend right=50, rr, "" name=D, "G_2 F_2"']
		& & \arrow[Rightarrow, to path=(U) -- (D) \tikztonodes, "\psi \phi"] \mathcal{E} .
	\end{tikzcd} \]
	Karena ini memberikan funktor, \eqref{eqn:dgCat-fct-compose} niscaya
	\begin{inparaenum}[(a)]
		\item asosiatif,
		\item komutatif dengan penarikan balik (atau pendorongan maju) sepanjang
		morfisme berderajat dalam kategori-dg
		$\iHom(\mathcal{D},\mathcal{E})$ (atau
		$\iHom(\mathcal{C},\mathcal{D})$).
	\end{inparaenum}
	Namun, operasi dalam (b) tidak lain adalah komposisi vertikal yang
	diperkenalkan dalam \eqref{eqn:dgfct-iHom}. Dengan demikian, kita melihat
	bahwa urutan komposisi vertikal dan horizontal dapat dipertukarkan.

	Untuk funktor antara kategori biasa, morfisme antara funktor-funktor itu juga memiliki
	komposisi vertikal dan horizontal yang saling komutatif. Hal ini dapat
	dipahami sebagai sifat kategori tertutup Kartesius $\cate{Cat}$ (atau
	dengan memandangnya sebagai kategori-$2$; lihat
	\cite[Contoh 3.5.3]{Li1}). Sifat $\cate{dgCat}_{\Bbbk}$ di atas adalah versi
	dg, atau versi aljabar linear, dari sifat tersebut.
\end{remark}
